\documentclass[12pt,a4paper]{article}

\usepackage[margin=1in]{geometry}
\usepackage{graphicx}
\usepackage{booktabs}
\usepackage{hyperref}
\usepackage[backend=biber,style=ieee]{biblatex}

\addbibresource{references.bib}

\title{Performance Isolation and Tail-Latency Analysis in Multi-Tenant Sidecarless Service Mesh Architectures}
\author{Author Name Placeholder\\
Department of Computer Science\\
University Placeholder}
\date{\today}

\begin{document}

\maketitle

\begin{abstract}
The contemporary transition from sidecar-based service meshes to sidecarless, kernel-accelerated designs promises lower per-request overhead, improved node density, and simplified data-plane management. Yet this efficiency shift introduces an unresolved systems tension: the same architectural centralization that reduces the sidecar tax may also increase the risk of cross-tenant performance interference when packet processing and policy enforcement converge in shared node-level components. This proposal investigates that tension through the lens of tail-latency isolation in multi-tenant environments, with particular emphasis on p99 behavior under adversarial and bursty workload conditions. The study focuses on eBPF/XDP-enabled service networking and sidecarless control patterns, including Cilium and Istio Ambient Mesh, to evaluate whether shared kernel-level data planes can deliver both operational efficiency and private performance guarantees. The intended contribution is an empirically grounded methodology and evidence base for assessing isolation quality, noisy-neighbor susceptibility, and mitigation strategies in next-generation cloud-native networking stacks.
\end{abstract}

\section{Introduction}
The service mesh has become a foundational mechanism for traffic management, observability, and policy enforcement in microservice deployments, but its original sidecar-centric realization has imposed measurable compute and latency overhead on production systems. This overhead, often described as the sidecar tax, emerges from duplicated proxies, additional context switches, memory footprints per pod, and increased request-path complexity under high fan-out communication graphs \cite{envoyDataPlaneOverhead,serviceMeshPerfStudies}. In multi-tenant clusters where economic efficiency and strict performance isolation must coexist, the sidecar tax is not merely a cost issue; it is an architectural tax that constrains packing density and distorts capacity planning.

Sidecarless approaches seek to reduce this tax by consolidating data-plane functionality into shared node-level agents and kernel-resident packet paths, frequently leveraging eBPF and XDP for low-overhead traffic processing. While this shift can reduce average latency, tail behavior remains the dominant concern for user-facing reliability. In this proposal, p99 tail-latency is adopted as the primary performance metric because it captures the rare but consequential slow requests that directly affect service-level objectives and user-perceived quality. The central argument is that mean latency improvements are insufficient if p99 latency exhibits inter-tenant leakage under contention; therefore, any claim of sidecarless superiority must be validated against isolation quality at the tail.

\section{Preliminary Literature Review}
Kernel-bypass and in-kernel fast-path research has established a strong foundation for contemporary cloud-native networking. The XDP model introduced by H\o iland-J\o rgensen \textit{et al.} demonstrated that programmable packet processing at early ingress can substantially reduce processing overhead and improve throughput under realistic traffic loads \cite{xdp2018}. This line of work, combined with sustained progress in eBPF verifier capabilities and tooling, has enabled production-grade policy enforcement and observability in systems such as Cilium \cite{ciliumArchitecture,ciliumHubble}.

In parallel, service mesh design has shifted from per-pod proxies toward shared or waypoint-based mediation models, including Istio Ambient Mesh and ztunnel-based transport indirection \cite{istioAmbient,ambientDesign}. These architectures reduce per-workload proxy replication and alter the reconciliation loop between control-plane intent and data-plane realization. Existing performance discussions frequently emphasize CPU and memory savings, startup improvements, and operational ergonomics, but comparatively fewer studies provide rigorous multi-tenant fairness analyses under controlled noisy-neighbor conditions. The literature therefore suggests a maturing implementation ecosystem but an incomplete evidence base regarding isolation guarantees in shared kernel-level planes.

\section{Research Gap}
The identified research gap is the lack of rigorous, reproducible, and tenant-aware empirical data on resource fairness and tail-latency isolation in sidecarless meshes that depend on shared kernel-level packet processing. Although prior work has established that eBPF/XDP can be fast, and platform guidance has highlighted the operational advantages of sidecarless patterns, there remains insufficient quantitative characterization of whether these gains persist when adversarial or bursty tenant behavior competes for common node-level agents, queues, and CPU cycles. Specifically, the field lacks systematic p99-focused comparison across sidecar-based and sidecarless configurations under equivalent traffic and resource pressure, along with kernel-level observability that can explain causality rather than merely report symptoms.

\section{Research Questions}
This thesis is guided by the primary question: to what extent do sidecarless service mesh architectures preserve p99 tail-latency isolation among tenants under shared-node contention when compared with sidecar-based baselines?

The investigation is further structured by three sub-questions. First, how does noisy-neighbor intensity alter p99 latency distribution for otherwise isolated namespaces under sidecarless data planes? Second, which kernel-level signals (for example queueing, drops, scheduling delay, and map update dynamics) most strongly correlate with observed latency leakage? Third, how does control-plane discovery freshness, including propagation lag in reconciliation loops, influence transient and steady-state tail-latency behavior during policy and endpoint churn events?

\section{Aims and Objectives}
The overarching aim is to produce a defensible, measurement-driven assessment of performance isolation in multi-tenant sidecarless meshes and to derive practical optimization guidance for cloud-native operators.

Objective 1 is to establish a baseline benchmarking suite that compares sidecar-based Envoy paths against sidecarless eBPF-based paths using identical traffic models, deployment topologies, and p99-centric evaluation criteria.

Objective 2 is to design and implement a configurable noisy-neighbor experimentation framework that can inject controlled CPU, memory, and traffic disturbances at tenant and node scopes while preserving repeatable test conditions.

Objective 3 is to perform kernel-level profiling using eBPF instrumentation, bpftrace probes, and mesh telemetry to map p99 degradations to specific low-level mechanisms in packet handling, scheduling, and policy enforcement.

Objective 4 is to formulate and evaluate an optimization proposal, potentially combining scheduling constraints, queueing controls, and policy tuning, with the goal of improving tail-latency isolation without reintroducing prohibitive architectural overhead.

\section{Research Approach (Methodology)}
The empirical study will employ an experimental systems methodology based on controlled, repeatable cluster deployments. A Kubernetes-in-Docker (Kind) environment will provide the primary testbed to ensure reproducibility and rapid scenario iteration, while still exposing realistic orchestration behavior for service discovery, policy updates, and workload scaling. Two principal architecture variants will be instantiated: a sidecar-based baseline and a sidecarless configuration using eBPF-enabled service networking with Cilium and Istio Ambient Mesh components, including ztunnel where applicable \cite{kindProject,istioAmbient,ciliumArchitecture}.

Load generation will use Fortio and wrk2 to capture both closed-loop and open-loop traffic regimes, enabling robust analysis of queue buildup and tail amplification under varying offered load \cite{fortio,wrk2}. Experimental factors will include tenant count, request concurrency, burstiness, resource saturation level, and control-plane update frequency. The primary response variable will be p99 latency, supplemented by p95, throughput, loss, and jitter to support interpretation. The design will prioritize repeated trials and confidence-aware reporting to reduce sensitivity to run-to-run variance.

Observation and diagnosis will combine application-level metrics with kernel-level evidence. bpftrace will be used for targeted tracepoint and kprobe-level inspection of scheduling and networking events, while Cilium Hubble will provide flow-level context for policy and path decisions \cite{bpftrace,ciliumHubble}. This dual-observability approach is intended to bridge macro-level latency behavior with micro-level causal signals. Statistical analysis will compare distributions across architectural modes and disturbance regimes, with particular attention to tail leakage patterns and the conditions under which sidecarless efficiency gains degrade into fairness trade-offs.

\section{Project Timeline}
Table~\ref{tab:timeline} outlines a 12-month execution plan aligned with engineering milestones and dissertation deliverables.

\begin{table}[h]
  \centering
  \caption{Proposed 12-month project timeline}
  \label{tab:timeline}
  \begin{tabular}{@{}ll@{}}
    \toprule
    \textbf{Months} & \textbf{Planned Activities} \\
    \midrule
    1--2 & Environment setup, tooling validation, baseline experiment design (Kind, mesh stacks, observability). \\
    3--5 & Development of noisy-neighbor framework, automation scripts, and reproducible scenario orchestration. \\
    6--8 & Large-scale experimental testing, repeated trial execution, and kernel-level trace collection. \\
    9 & Data cleaning, statistical analysis, and synthesis of p99 isolation findings. \\
    10--12 & Thesis writing, revision, final optimization evaluation, and defense preparation. \\
    \bottomrule
  \end{tabular}
\end{table}

\section{Expected Contribution}
The expected contribution is a technically rigorous and practically relevant account of how sidecarless service mesh architectures behave under multi-tenant contention, with emphasis on tail-latency isolation rather than average-path efficiency alone. By integrating controlled benchmarking with kernel-level profiling, the study aims to provide both empirical validation and mechanistic explanation. The resulting artifacts, including workload definitions and measurement procedures, are intended to support reproducible follow-on research and evidence-based operational guidance for performance-sensitive cloud-native platforms.

\printbibliography

\end{document}
